\documentclass[10pt,a4paper]{extarticle}
\usepackage{parskip}

\usepackage[utf8]{inputenc}
\usepackage[T1]{fontenc}
\usepackage[brazil]{babel}
\usepackage{amsmath, amssymb, amsthm}
\usepackage{geometry}
\usepackage{xcolor}
\usepackage{hyperref}
\usepackage{booktabs}
\usepackage{array}
\usepackage{tabularx}
\usepackage{enumitem}
\usepackage[protrusion=true,expansion=false]{microtype}
\usepackage{csquotes}
\usepackage{tcolorbox}
\usepackage{natbib}

\definecolor{important}{RGB}{255, 100, 100}
\definecolor{notice}{RGB}{0, 102, 204}
\definecolor{accent}{RGB}{230, 120, 30}

\newcommand{\impt}[1]{\textcolor{important}{#1}}
\newcommand{\ntc}[1]{\textcolor{notice}{#1}}

% Comando para instruções de preenchimento -- REMOVA todas as ocorrências
% de \guia{} antes de entregar o documento final.
\newcommand{\guia}[1]{{\small\ntc{\textit{[Orientação -- remover antes de entregar: #1]}}}}

\setlist{nosep}
\geometry{margin=2.5cm}

\newtcolorbox{respbox}{
  colback=notice!5,
  colframe=notice!50,
  boxrule=0.4pt,
  arc=2pt,
  left=4pt, right=4pt, top=3pt, bottom=3pt,
  fontupper=\small
}

\title{Proposta de Pesquisa\\ \large \guia{substitua pelo título do seu trabalho}}
\author{\guia{Nome completo do(a) aluno(a)} \\ BCC502 -- Metodologia Científica para Ciência da Computação}
\date{\today}

\begin{document}

\maketitle

\begin{center}
\ntc{\textit{Este documento é um template. Preencha cada seção conforme as instruções em cor azul,
e apague todas as instruções (\guia{...}) antes da entrega final. Com exceção da seção de
Objetivos, todo o texto deve ser redigido em prosa corrida -- não utilize enumerações ou
tópicos para responder.}}
\end{center}

\section{Contextualização e Problema de Pesquisa}

\guia{Texto corrido, entre 1 e 2 páginas, sem subdivisão em tópicos. Apresente, de forma
articulada, o domínio de aplicação ou área de pesquisa; o problema que motiva a investigação;
a relevância desse problema do ponto de vista científico, tecnológico ou social; pelo menos
três trabalhos relacionados que evidenciem o estado atual do conhecimento sobre o tema
(cite-os com \textbackslash citep\{\}); e as limitações, lacunas ou oportunidades de pesquisa
que ainda permanecem em aberto. O texto deve caminhar do domínio geral até a lacuna
específica que esta pesquisa pretende preencher, como argumentação única.}

[Escreva aqui o texto de contextualização.]

\section{Objetivos}

\ntc{\guia{Esta é a única seção do documento em que a resposta pode ser apresentada em
formato de lista.}}

\subsection{Objetivo Geral}

\guia{Uma única frase, iniciada com verbo no infinitivo, verificável ao final do trabalho.}

[Escreva aqui o objetivo geral.]

\subsection{Objetivos Específicos}

\guia{Entre três e cinco objetivos, em lista numerada. Cada um deve ser um desdobramento
mensurável do objetivo geral -- não descreva aqui etapas do método (ex.: "implementar o
protótipo" não é um objetivo específico).}

\begin{enumerate}[label=\textbf{OE\arabic*.}]
    \item{} [Objetivo específico 1.]
    \item{} [Objetivo específico 2.]
    \item{} [Objetivo específico 3.]
\end{enumerate}

\guia{Feche a seção com um parágrafo corrido explicando como os objetivos específicos,
em conjunto, contribuem para atingir o objetivo geral.}

[Escreva aqui o parágrafo de articulação entre objetivos específicos e objetivo geral.]

\section{Hipótese de Trabalho}

\guia{Texto corrido, sem subdivisão em tópicos. Apresente a hipótese como uma única frase
declarativa; justifique sua plausibilidade à luz da literatura consultada; explique como ela
poderá ser investigada ou testada; e discuta quais evidências seriam suficientes para
aceitá-la ou rejeitá-la.}

[Escreva aqui o texto sobre a hipótese de trabalho.]

\section{Estratégia de Investigação e Validação}

\guia{Texto corrido descrevendo a estratégia preliminar para investigar ou validar a
hipótese (dados, experimentos, protótipos, critérios de comparação) e, na sequência, as
limitações dessa estratégia -- o que ela não é capaz de garantir ou controlar.}

[Escreva aqui a estratégia de investigação/validação e suas limitações.]

\section{Coerência da Proposta}

\guia{Parágrafo de encerramento, em texto corrido, evidenciando explicitamente a coerência
entre o problema de pesquisa, o objetivo geral, os objetivos específicos, a hipótese de
trabalho e a estratégia de validação -- mostrando como cada elemento decorre logicamente do
anterior, e não apenas listando-os lado a lado.}

[Escreva aqui o parágrafo de coerência.]

\section{Utilização de IA}

\guia{Texto corrido relatando o uso de ferramentas de IA generativa na elaboração desta
proposta -- mesmo que você não tenha utilizado nenhuma, declare isso explicitamente. Deixe
claro o que foi pensado/redigido por você, o que foi produzido com apoio de IA, e como você
validou ou corrigiu o que a IA produziu antes de incorporá-lo ao texto final.}

[Escreva aqui o relato sobre o uso de IA.]

\guia{Em seguida, reproduza integralmente, na ordem cronológica de uso, todos os prompts
fornecidos à IA -- sem edição de conteúdo, incluindo eventuais erros de digitação. Use um
bloco \texttt{respbox} por prompt. Duplique o bloco abaixo quantas vezes forem necessárias.}

\begin{respbox}
\textbf{Prompt 1:} [cole aqui, sem alterações, o texto exato do primeiro prompt utilizado]
\end{respbox}

\begin{respbox}
\textbf{Prompt 2:} [cole aqui o segundo prompt]
\end{respbox}

%\bibliographystyle{apalike}
%\bibliography{references}

\end{document}